\chapter{System Design and Implementation}\label{ch:system}

\section{Overview}

Figure~\ref{fig:architecture} shows the system. The robot computer, a laptop, reads two cameras and the arm's joint
positions, sends observations to a GPU server, and executes the returned chunks of actions on the arm at 30~Hz. Between
the two, a fault injector can delay, drop, duplicate or reorder messages. Everything on the robot computer except the
scheduler and the smoothing is LeRobot~0.7 as it ships; the server is LeRobot's policy server, unchanged.

\begin{figure}[htbp]
  \centering
  % System overview: robot computer, network with fault injection, GPU server. Colour is allowed in figures.
\definecolor{robotfill}{HTML}{EAF2FC}
\definecolor{robotline}{HTML}{2A78D6}
\definecolor{serverfill}{HTML}{FDEEE7}
\definecolor{serverline}{HTML}{EB6834}
\definecolor{netfill}{HTML}{F0EFEC}
\definecolor{ink}{HTML}{0B0B0B}
\definecolor{ink2}{HTML}{52514E}
\begin{tikzpicture}[
  font=\small, node distance=4mm and 6mm,
  blk/.style={draw=#1, rounded corners=2pt, thick, fill=white, align=center, minimum height=8mm, inner sep=3pt},
  arr/.style={-{Latex[length=2.2mm]}, thick, ink2},
]
% robot computer
\node[blk=robotline, minimum width=27mm] (cams) {Cameras\\ wrist, scene};
\node[blk=robotline, minimum width=27mm, below=of cams] (arm) {SO-101 arm\\ (follower)};
\node[blk=robotline, minimum width=30mm, right=10mm of cams] (sched) {Scheduler\\ upstream \textbar{} \drtc};
\node[blk=robotline, minimum width=30mm, below=of sched] (queue) {Action queue};
\node[blk=robotline, minimum width=30mm, below=of queue] (smooth) {Smoothing\\ filter, playback};
\begin{scope}[on background layer]
  \node[draw=robotline, fill=robotfill, rounded corners=4pt, thick, fit=(cams)(arm)(sched)(queue)(smooth),
        inner sep=5mm, label={[font=\small\bfseries, text=ink]above:Robot computer (laptop), 30~Hz}] (rc) {};
\end{scope}
% network
\node[draw=ink2, fill=netfill, rounded corners=3pt, thick, minimum width=24mm, minimum height=33mm, align=center,
      right=14mm of queue] (net) {Network\\ (WiFi)\\[6pt] \textbf{fault}\\ \textbf{injection}};
% server
\node[blk=serverline, minimum width=30mm, right=14mm of net, yshift=8mm] (ps) {Policy server\\ (LeRobot 0.7)};
\node[blk=serverline, minimum width=30mm, below=of ps] (pol) {Policy\\ \smolvla{} \textbar{} \groot};
\begin{scope}[on background layer]
  \node[draw=serverline, fill=serverfill, rounded corners=4pt, thick, fit=(ps)(pol), inner sep=5mm,
        label={[font=\small\bfseries, text=ink]above:GPU server (A100)}] (srv) {};
\end{scope}
% flows
\draw[arr] (cams) -- node[above, font=\scriptsize]{images} (sched);
\draw[arr] (sched.east) -- ++(4mm,0) |- ([yshift=4mm]net.west) node[pos=0.75, above, font=\scriptsize]{observation};
\draw[arr] ([yshift=4mm]net.east) -- node[above, font=\scriptsize]{} ([yshift=-2mm]ps.west);
\draw[arr] (ps) -- (pol);
\draw[arr] ([yshift=-4mm]pol.west) -- ++(-4mm,0) |- ([yshift=-6mm]net.east);
\draw[arr] ([yshift=-6mm]net.west) -- node[below, font=\scriptsize, pos=0.45]{chunk ($H$ actions)} (queue.east);
\draw[arr] (queue) -- (smooth);
\draw[arr] (smooth.west) -| node[pos=0.25, below, font=\scriptsize]{commands} (arm.south);
\draw[arr, densely dotted] (arm.north) -- node[left, font=\scriptsize]{joint state} (cams.south);
\end{tikzpicture}

  \caption{The system. The robot computer runs the control loop at 30~Hz; the scheduler decides when to request the
           next chunk, and the smoothing shapes the commands sent to the arm. Faults are injected where a real network
           would cause them, so both the shipped client and \drtc{} react to them with their own logic.}
  \label{fig:architecture}
\end{figure}

\section{Hardware}\label{sec:hardware}

\begin{table}[htbp]
  \centering
  \caption{Hardware and software.}
  \label{tab:hardware}
  \begin{tabularx}{\textwidth}{l X}
    \toprule
    Component & Description \\
    \midrule
    Arm & SO-101 follower arm~\cite{so101} with six Feetech servos; an identical leader arm for recording \\
    Cameras & Wrist: Zebronics Zeb Live Pro, $800\times600$ MJPG at 30~fps. Scene: an Android phone over USB
              (DroidCam), $800\times600$ at 30~fps \\
    Robot computer & Laptop (Linux, NVIDIA RTX~4060 Laptop GPU, 8~GB), running the control loop and the client;
                     used for every experiment in this report \\
    Model server & NVIDIA A100 80~GB (a shared GPU server) for every robot experiment \\
    Training and offline tests & NVIDIA H200 (141~GB, shared, in a container) and the A100 \\
    Network & Laptop to server over WiFi and the local network. Each request carries both camera images
              uncompressed ($2\times800\times600\times3$~bytes $\approx$ 2.9~MB) and the joint state \\
    Software & LeRobot~0.7.0, PyTorch~2.14.1, Python~3.12; our layer, the \texttt{lab} command (Section~\ref{sec:impl}) \\
    \bottomrule
  \end{tabularx}
\end{table}

\begin{figure}[htbp]
  \centering
  \photo[0.9\textwidth]{setup.jpg}{The whole setup: the SO-101 arm, the task board, the cameras and the laptop.}
  \caption{The experimental setup.}
  \label{fig:setup}
\end{figure}

\begin{figure}[htbp]
  \centering
  \begin{subfigure}[t]{0.48\textwidth}
    \centering
    \photo[\textwidth]{arm_wrist_camera.jpg}{The follower arm with the wrist camera.}
    \caption{Follower arm and wrist camera.}
  \end{subfigure}\hfill
  \begin{subfigure}[t]{0.48\textwidth}
    \centering
    \photo[\textwidth]{scene_camera.jpg}{The scene camera (phone) and its view of the task.}
    \caption{Scene camera.}
  \end{subfigure}
  \caption{The robot and its cameras.}
  \label{fig:robot}
\end{figure}

\section{The baseline: LeRobot remote inference}\label{sec:baseline}

The baseline is LeRobot~0.7's remote-inference client as it ships. With $q_t$ queued actions, an action period
$\Delta = 1/30$~s and a turnaround estimate $\hat L$, it requests the next chunk when
\begin{equation}
  q_t\,\Delta \;\le\; \max\!\left(T_{\mathrm{refill}},\; \hat L + \Delta\right),
\end{equation}
where $T_{\mathrm{refill}}=0.5$~s and $\hat L$ is the largest of the last 100 measured turnarounds. It keeps one request
in flight; a reply that misses its deadline is a fault; and when no action is available for the starvation grace of
1~s, it ends the run.

\section{\drtc{} on LeRobot 0.7}\label{sec:drtc}

We implemented \drtc's scheduling inside LeRobot's own runtime, replacing only the decision of when to request and
when to give up on a request. All of LeRobot's checks stay in place: one request in flight, observation freshness,
generations, deadlines and latched faults.

\paragraph{Trigger.} A chunk is requested when $q_t \le H - s_{\min}$, with $H$ the length of the last accepted chunk.
We use $s_{\min}=20$ steps for \smolvla{} ($H=50$) and $s_{\min}=8$ for \groot{} ($H=16$).

\paragraph{Latency estimate.} For every measured turnaround $R$, the Jacobson--Karels estimator updates a smoothed mean
$\bar R$ and a smoothed mean deviation $D$,
\begin{equation}
  D \leftarrow (1-\beta)\,D + \beta\,|R-\bar R|, \qquad \bar R \leftarrow (1-\alpha)\,\bar R + \alpha R, \qquad
  \hat L = \bar R + k\,D,
\end{equation}
with $\alpha=1/8$, $\beta=1/4$ (as in TCP~\cite{rfc6298}) and $k=1.5$.

\paragraph{Cooldown.} A request not answered within its cooldown $c$ is given up and a fresh observation is sent. In
\drtc{} the cooldown is the latency estimate in steps, capped at $H/2$, plus a margin of $\varepsilon$ steps. Ours is
\begin{equation}
  c = \max\!\left(2,\; \left\lceil \tilde L / \Delta \right\rceil + \varepsilon\right)\Delta, \qquad
  \tilde L = \max\!\left(\hat L,\; \max_{i\,\in\,\text{last 10}} R_i\right),
\end{equation}
with $\varepsilon = 1$, and before the first reply $c = \max\!\left((s_{\min}+\varepsilon)\Delta,\;
T_{\mathrm{start}}\right)$ with the start-up timeout $T_{\mathrm{start}}=10$~s.

\paragraph{Holding through stalls.} Under \drtc{} the arm holds its last commanded position while no action is
available and continues when a chunk arrives; the run is not ended after one second.

\paragraph{Fixes found on the real setup.} Table~\ref{tab:fixes} lists the four changes to \drtc{} that the
experiments required. Each was found as a failure on the real robot and verified on the replay bench before the main
experiment.

\begin{table}[htbp]
  \centering
  \caption{Failure modes of \drtc{} on the real setup and our fixes.}
  \label{tab:fixes}
  \begin{tabularx}{\textwidth}{>{\raggedright\arraybackslash}p{0.21\textwidth} X X}
    \toprule
    Failure & Cause & Fix \\
    \midrule
    \groot{} never received a chunk & The cooldown is capped at $H/2$: 8 steps (0.27~s) for a 16-action chunk, while
      replies took 0.4--0.7~s, so every request was given up before its reply & The cooldown is never shorter than the
      measured latency (no cap; the original cap remains an option) \\
    A storm of re-requests & Replies close to the estimate were given up; the server, still computing them, slowed
      further & The cooldown also covers the slowest of the last 10 replies \\
    The robot never started on a slow link & Before the first reply the cooldown is only $s_{\min}+\varepsilon$ steps
      & Until the first reply, wait up to the start-up timeout \\
    \enquote{Invalid RTC delay} errors & The requested inference delay could exceed the server's chunk & Cap the
      delay at the server's prediction steps \\
    \bottomrule
  \end{tabularx}
\end{table}

\section{Smoothing}\label{sec:smoothing}

\paragraph{Butterworth filter.} Every commanded arm position passes through a causal second-order Butterworth low-pass
filter per joint, with a cutoff of 3~Hz at the 30~Hz control rate, implemented as second-order sections. A gain $g=1.4$
on each step compensates the filter's loss of amplitude: with $y_t$ the filter output, the command is
$u_t = u_{t-1} + g\,(y_t - u_{t-1})$. The filter keeps its state across chunk boundaries, which smooths the switch from
one chunk to the next, and is reset when motion is invalidated.

\paragraph{Adaptive playback.} When the queue would empty before the next chunk arrives, the actions are played more
slowly instead of the arm stopping. With $n_t$ the expected ticks until the next chunk (from the latency estimate and
the time since the request was sent), the playback rate is
\begin{equation}
  r^\star_t = \mathrm{clip}\!\left(\frac{q_t}{n_t + m},\; r_{\min},\; 1\right), \qquad
  r_t = r_{t-1} + \mathrm{clip}\!\left(r^\star_t - r_{t-1},\, -\sigma,\, \sigma\right),
\end{equation}
with $r_{\min}=0.5$, a margin $m=4$ steps and a slew limit $\sigma=0.05$ per tick. A phase $\phi$ advances by $r_t$
per tick, and the command is interpolated linearly between the two actions around it. With a full queue, $r_t=1$ and
playback is unchanged.

\section{Network fault injection}\label{sec:faults}

Faults are injected in the client's transport, where a real network would cause them, so that the scheduler, the
deadlines and the starvation logic all respond as they would to a real fault. The injector supports dropping,
duplicating and reordering observations and replies, latency spikes (a one-time extra delay of a reply), disconnects
(both directions dropped) and, on the bench, a constant latency with jitter. Every injected fault is logged.

The experiments use one fixed \emph{fault mix} (Figure~\ref{fig:faultmix}), on the trial's clock: a latency spike of
+1~s at 2~s, every reply dropped from 4 to 5~s, and a full outage from 6 to 7.5~s.

\begin{figure}[htbp]
  \centering
  % The fault mix of conditions 2-4 on the trial's clock (seconds from the trial start).
\definecolor{band}{HTML}{E1E0D9}
\definecolor{ink}{HTML}{0B0B0B}
\definecolor{ink2}{HTML}{52514E}
\definecolor{spike}{HTML}{EB6834}
\begin{tikzpicture}[x=11mm, y=7mm, font=\small]
  \draw[ink2, thick, -{Latex[length=2mm]}] (0,0) -- (10.6,0) node[right]{s};
  \foreach \t in {0,...,10} { \draw[ink2] (\t,0) -- (\t,-0.15) node[below, font=\scriptsize]{\t}; }
  % latency spike: the reply due after 2 s arrives 1 s late
  \draw[spike, very thick, -{Latex[length=2mm]}] (2,0.35) -- (3,0.35);
  \fill[spike] (2,0.35) circle (2pt);
  \node[above, align=center, font=\scriptsize] at (2.5,0.5) {latency spike\\ +1~s at 2~s};
  % every reply dropped 4-5 s
  \fill[band, draw=ink2] (4,0.15) rectangle (5,0.85);
  \node[above, align=center, font=\scriptsize] at (4.5,0.9) {replies dropped\\ 4--5~s};
  % full outage 6-7.5 s
  \fill[band, draw=ink2] (6,0.15) rectangle (7.5,0.85);
  \draw[ink2] (6,0.15) -- (7.5,0.85) (6,0.85) -- (7.5,0.15);
  \node[above, align=center, font=\scriptsize] at (6.75,0.9) {outage, both directions\\ 6--7.5~s};
\end{tikzpicture}

  \caption{The fault mix of conditions 2--4, in seconds from the start of a trial.}
  \label{fig:faultmix}
\end{figure}

\section{Evaluation framework}\label{sec:evalsys}

\paragraph{Evaluation plans.} Each evaluation is a file that fixes the task, the written success rule, the time limit,
the policies and the conditions. Before the first trial it is \emph{frozen}: every policy is pinned to its exact model
version, so all trials of an evaluation use the same models.

\paragraph{Rounds.} Trials run in rounds. Each round runs every policy under every condition once, in a seeded random
order, so slow drifts (light, motor temperature) affect all conditions alike. The operator resets the scene, starts
the trial, and presses a key for success or failure.

\paragraph{Records.} Every trial stores its raw logs (one row per control tick and per chunk, the injected faults, the
scheduler's recovery actions), its configuration and a video from each camera. Raw files are never edited; a trial
with a setup mistake is marked void with a reason, kept, and listed in the report. Sessions are pushed to a shared,
private results store, and the reports are rebuilt from the raw files, so every number in this report can be traced to
its trials.

\paragraph{Bench.} A replay robot plays recorded held-out episodes as the camera input, without an arm, so the network
behaviour of every condition can be measured without operator time. It reports the timing and smoothness metrics but
not task success.

\paragraph{Offline evaluation.} Without a robot, each trained policy predicts chunks from held-out recorded
observations; the predictions are compared with the recorded actions, and the inference time and GPU memory are
measured.

\section{Implementation}\label{sec:impl}

The project is a Python layer on top of LeRobot~0.7, which is used as a pinned dependency. A registry of YAML files
describes the datasets, training recipes, models, rigs, machines, experiment conditions and evaluation plans; a single
command, \texttt{lab}, trains, serves, runs and evaluates from it. The \drtc{} engine, the fault injector, the smoothing
and the evaluation framework are covered by automated tests. Appendix~\ref{app:repro} lists the commands that reproduce
the results.
